\scanpage{66}
d) \deleted{\unsure{a}} Les $\Sigma_j\hookrightarrow\Sigma_{ji}$ sont bordants, et si
$x\in\Sigma_i$, \deleted{les} donc les $\Sigma_{ji}\ni x$ (\deleted{pour $j<i$}\inserted{indices}
\deleted{tels que}\deleted{$\unsure{\Sigma}$}\inserted{par l'ens} $I_{<i,x}\overset{\mathrm{def}}{=}\{j\in I\mid j<i,\ \Sigma_{ji}\ni x\}$,
qui est totalement ordonné par b)) sont
des sous-espaces-bords transversaux en $x$
--- i.e. les $\Sigma_{ji}$ $(j\in I_{<i})$ sont des
sous-espaces bords de $\Sigma_i$ mutuellement
transversaux.

e) $\Sigma_{ij}\xrightarrow{s_{ij}}\Sigma_i$ est transversal p.v. au système trans-
versal des morphismes bordants
$\unsure{b_{ki}|b_{ik}}:\unsure{\Sigma_{ki}}\longrightarrow\Sigma_i$

Si on pose
\[
\Sigma=\coprod_{i\in I}\Sigma_i
\]
\[
\Sigma^{<}(1)=\coprod_{\substack{i,j\in I\\i<j}}\Sigma_{ij}
\]
et $\uline{s},\uline{b}:\Sigma^{<}(1)\longrightarrow\Sigma$ induits par les
$s_{ij},b_{ij}$, alors les conditions a) b) c)
signifient que $(\uline{s},\uline{b}):\Sigma^{<}(1)\longrightarrow\Sigma\times\Sigma$ est
inj. et a comme \deleted{graphe}\inserted{image} le graphe d'une
relation d'ordre stricte, avec $\uline{b}:\Sigma^{<}(1)\longrightarrow\Sigma$
\inserted{immersion} propre, \deleted{et} que l'application \inserted{continue évidente} $\varphi:\unsure{\Sigma|\Sigma^{*}}\longrightarrow I$
est strict${}^t$ croissante, et que pour tt $x\in\Sigma$,
l'application induite $\varphi:\Sigma_{<x}\longrightarrow I_{<i}$ $(i=\varphi(x))$ est
injective, enfin que $\Sigma_{<x}$ est tot${}^t$ ordonné.
\mspage{65}
